The analysis considers a uniform tube of length $L$ and constant cross-sectional area $A$ closed rigidly at both ends ($x=0$ and $x=L$), as shown in Figure~\ref{fig:hw02-03}. The tube contains an inviscid fluid of equilibrium density $\rho$, bulk modulus $K$, and acoustic sound speed $c=\sqrt{K/\rho}$. Small-amplitude axial vibrations $u(x,t)$ satisfy the one-dimensional wave equation established in Problems B.1 and B.2.

\textbf{1. Boundary conditions in terms of displacement.} At a rigid, impermeable wall, fluid particles cannot move into or out of the boundary. Therefore, the axial displacement must vanish at both closed ends for all time:
\[
\boxed{u(0,t)=0}\qquad\text{(left closed end, B.3.1.a)}
\]
\[
\boxed{u(L,t)=0}\qquad\text{(right closed end, B.3.1.b)}
\]

\textbf{2. Harmonic solution of the wave equation.} From Problem B.2 (item 14), the monochromatic standing-wave solution consists of the product of a spatial function $U(x)$ and a temporal harmonic factor $T(t)$:

\textbf{(a) First term (spatial distribution):}
\[
\boxed{U(x)=C_1\sin(\gamma x)+C_2\cos(\gamma x)}\qquad\text{(B.3.2.a)}
\]
where $\gamma$ is the wavenumber and $C_1, C_2$ are arbitrary constants.

\textbf{(b) Second term (temporal harmonic oscillation):}
\[
\boxed{T(t)=e^{-i\omega t}}\qquad\text{(B.3.2.b)}
\]
where $\omega$ is the angular frequency.

The complete displacement solution is therefore
\[
\boxed{u(x,t)=U(x)T(t)=\bigl(C_1\sin\gamma x+C_2\cos\gamma x\bigr)e^{-i\omega t}}\qquad\text{(B.3.2.c)}
\]

\textbf{3. Boundary conditions in terms of constants $C_1$ and $C_2$.} Substituting the harmonic solution into the boundary conditions at $x=0$ and $x=L$ gives:
\[
u(0,t)=\bigl(C_1\sin 0+C_2\cos 0\bigr)e^{-i\omega t}=C_2e^{-i\omega t}=0\qquad\text{(B.3.3.a)}
\]
\[
u(L,t)=\bigl(C_1\sin\gamma L+C_2\cos\gamma L\bigr)e^{-i\omega t}=0\qquad\text{(B.3.3.b)}
\]

\textbf{4. Determination of $C_2$ and equation for $C_1$.} Because $e^{-i\omega t}\neq 0$ for all $t$, condition (B.3.3.a) requires
\[
\boxed{C_2=0}\qquad\text{(B.3.4.a)}
\]
Substituting $C_2=0$ into condition (B.3.3.b) eliminates the cosine term and yields
\[
\boxed{C_1\sin(\gamma L)=0}\qquad\text{(B.3.4.b)}
\]

\textbf{5. Discussion of the trivial solution.} If $C_1=0$, then since $C_2=0$, the spatial function vanishes identically: $U(x)\equiv 0$, and consequently $u(x,t)\equiv 0$ for all $x\in[0,L]$ and all $t$. This trivial solution represents the fluid at rest in equilibrium, in which no vibration, acoustic disturbance, or sound wave is present.

\textbf{6. Nontrivial solution and governing condition.} For a sustained, physically meaningful vibration to exist, we seek a nontrivial solution where the amplitude is nonzero, $C_1\neq 0$. Dividing (B.3.4.b) by $C_1$ requires that the wavenumber $\gamma$ satisfy
\[
\boxed{\sin(\gamma L)=0}\qquad\text{(B.3.6)}
\]
When this condition holds, $C_1$ can take any arbitrary nonzero value. This reflects the single degree of indeterminacy characteristic of linear modal analysis, where absolute vibration amplitudes are set by external excitation or initial conditions rather than boundary equations alone.

\textbf{7. Eigenvalue equation.} Introduce the dimensionless parameter $z$ defined by
\[
\boxed{z=\gamma L}\qquad\text{(B.3.7.a)}
\]
Substituting this definition into the condition from item 6 yields the eigenvalue equation:
\[
\boxed{\sin z=0}\qquad\text{(B.3.7.b)}
\]

\textbf{8. Eigenvalue solutions.} The mathematical roots of $\sin z=0$ are $z=n\pi$ for any integer $n$. For physical vibrations characterized by positive propagation wavenumber ($\gamma>0$) and positive length ($L>0$), we restrict attention to the positive integers:
\[
\boxed{z_j=j\pi,\qquad j=1,2,3,\ldots}\qquad\text{(B.3.8)}
\]
The values $z_1=\pi, z_2=2\pi, z_3=3\pi, \ldots$ are the eigenvalues of the closed-closed tube system.

\textbf{9. General expression for the wavenumbers.} Using the relation $z_j=\gamma_j L$, the modal wavenumbers are
\[
\boxed{\gamma_j=\frac{z_j}{L}=\frac{j\pi}{L},\qquad j=1,2,3,\ldots}\qquad\text{(B.3.9)}
\]

\textbf{10. Definition of wavenumber.} In acoustics and wave mechanics (HW01 and Problem B.2), the wavenumber $\gamma$ relates the circular frequency $\omega$ to the sound speed $c$ via
\[
\boxed{\gamma=\frac{\omega}{c}}\qquad\text{(B.3.10)}
\]

\textbf{11. Angular frequency in terms of wavenumber and wavespeed.} Solving (B.3.10) for $\omega$ gives
\[
\boxed{\omega=\gamma c}\qquad\text{(B.3.11)}
\]

\textbf{12. General expression for the natural frequencies.} Substituting the modal wavenumbers $\gamma_j$ from item 9 into the frequency relation gives the natural angular frequencies (in $\mathrm{rad/s}$):
\[
\boxed{\omega_j=\gamma_j c=\frac{j\pi c}{L},\qquad j=1,2,3,\ldots}\qquad\text{(B.3.12)}
\]

\textbf{13. General expression for cyclic frequencies in Hz.} Because $\omega=2\pi f$, the cyclic frequency in hertz ($\mathrm{Hz}$ or $\mathrm{s}^{-1}$) is $f=\omega/(2\pi)$. Substituting (B.3.12) yields
\[
\boxed{f_j=\frac{\omega_j}{2\pi}=\frac{jc}{2L},\qquad j=1,2,3,\ldots}\qquad\text{(B.3.13)}
\]

\textbf{14. Meaning of fundamental frequency and overtones.}
\par\smallskip\noindent
\textbf{Fundamental frequency:} The lowest resonant frequency of the system, obtained for the lowest modal index $j=1$:
\[
f_1=\frac{c}{2L}.
\]
It represents the principal natural pitch produced by the vibrating air column.
\par\smallskip\noindent
\textbf{Overtones:} Any resonant frequency higher than the fundamental frequency ($j=2,3,4,\ldots$). The first overtone is $f_2=2f_1=c/L$; the second overtone is $f_3=3f_1=3c/(2L)$; and in general, the $n$-th overtone corresponds to index $j=n+1$ with frequency $f_{n+1}=(n+1)f_1$.

\textbf{15. Meaning of fundamental harmonic and higher harmonics.}
\par\smallskip\noindent
\textbf{Fundamental harmonic:} A harmonic is a resonance whose frequency is an integer multiple of the fundamental frequency $f_1$. The fundamental harmonic (also called the first harmonic) is $f_1=1\cdot f_1$.
\par\smallskip\noindent
\textbf{Higher harmonics:} Resonance frequencies that are integer multiples of the fundamental: the second harmonic is $2f_1$, the third harmonic is $3f_1$, etc.
\par\smallskip\noindent
\textbf{Comparison for closed-closed tube:} Because all positive integer multiples $j=1,2,3,\ldots$ satisfy the eigenvalue condition for a tube closed at both ends, all harmonics are present in the spectrum. Thus, the 1st overtone ($j=2$) is the 2nd harmonic, the 2nd overtone ($j=3$) is the 3rd harmonic, and the $n$-th overtone is the $(n+1)$-th harmonic.

\textbf{16. Relationship between wavelength and wavenumber.} The acoustic wavelength $\lambda$ and wavenumber $\gamma$ are connected by
\[
\boxed{\lambda=\frac{2\pi}{\gamma}}\qquad\text{(B.3.16)}
\]

\textbf{17. General expression for modal wavelengths.} Substituting the wavenumber $\gamma_j=j\pi/L$ from item 9 into (B.3.16) gives
\[
\boxed{\lambda_j=\frac{2\pi}{\gamma_j}=\frac{2\pi}{j\pi/L}=\frac{2L}{j},\qquad j=1,2,3,\ldots}\qquad\text{(B.3.17)}
\]
Higher modes have progressively shorter wavelengths.

\textbf{18. Wavelength multiples accommodating standing waves.} Rearranging (B.3.17) in terms of the tube length $L$ yields
\[
\boxed{L=j\left(\frac{\lambda_j}{2}\right),\qquad j=1,2,3,\ldots}\qquad\text{(B.3.18)}
\]
This relation shows that sustained vibratory standing waves can exist inside a tube closed at both ends if and only if the length of the tube accommodates an exact integer number of half-wavelengths:
\[
L=\frac{\lambda_1}{2}=2\left(\frac{\lambda_2}{2}\right)=3\left(\frac{\lambda_3}{2}\right)=\dots=j\left(\frac{\lambda_j}{2}\right).
\]
Because both boundaries are rigid walls where displacement must be zero (displacement nodes), every valid standing wave must place a displacement node at each end, forcing the distance between the ends to be an integral multiple of node-to-node spacing ($\lambda/2$).

\textbf{19. Elimination of $C_2$ from spatial form.} Recalling $U(x)=C_1\sin\gamma x+C_2\cos\gamma x$ from item 2(a) and inserting $C_2=0$ from item 4 gives
\[
\boxed{U(x)=C_1\sin(\gamma x)}\qquad\text{(B.3.19)}
\]

\textbf{20. Displacement modeshapes.} Setting the arbitrary scaling factor $C_1=1$ and evaluating at wavenumber $\gamma_j=j\pi/L$, the general expression for the $j$-th displacement modeshape is
\[
\boxed{U_j(x)=\sin(\gamma_j x)=\sin\left(\frac{j\pi x}{L}\right),\qquad j=1,2,3,\ldots}\qquad\text{(B.3.20)}
\]
At both ends, $U_j(0)=\sin(0)=0$ and $U_j(L)=\sin(j\pi)=0$, strictly fulfilling the closed-end boundary conditions.

\textbf{21. General expression of the pressure modeshapes.} From fluid elasticity (Problem B.1, Eq. B.1.5), the acoustic pressure disturbance $p(x,t)$ is related to spatial strain by
\[
p(x,t)=-Ku_x(x,t)=-\rho c^2 u_x(x,t).
\]
For harmonic mode $j$, substituting $u_j(x,t)=U_j(x)e^{-i\omega_j t}$ gives
\[
p_j(x,t)=P_j(x)e^{-i\omega_j t},\qquad\text{where}\quad P_j(x)=-\rho c^2 U_j'(x).
\]
Differentiating the displacement modeshape $U_j(x)=\sin(\gamma_j x)$ yields $U_j'(x)=\gamma_j\cos(\gamma_j x)$. Therefore, the physical pressure modeshape is
\[
\boxed{P_j(x)=-\rho c^2 \gamma_j\cos(\gamma_j x)=-\rho c\,\omega_j\cos\left(\frac{j\pi x}{L}\right)}\qquad\text{(B.3.21.a)}
\]
When scaled to unit amplitude (or written as a normalized pressure distribution $\bar{P}_j(x)=P_j(x)/|P_j|_{\max}$), the pressure modeshape simplifies to
\[
\boxed{\bar{P}_j(x)=-\cos\left(\frac{j\pi x}{L}\right)\quad\text{or}\quad \cos\left(\frac{j\pi x}{L}\right)}\qquad\text{(B.3.21.b)}
\]
This result reveals the fundamental spatial phase shift of $\pi/2$ between displacement and pressure: the closed ends ($x=0, L$) are displacement nodes ($U_j=0$) but pressure antinodes ($|\bar{P}_j|=1$).
